% Options for packages loaded elsewhere
% Options for packages loaded elsewhere
\PassOptionsToPackage{unicode}{hyperref}
\PassOptionsToPackage{hyphens}{url}
\PassOptionsToPackage{dvipsnames,svgnames,x11names}{xcolor}
%
\documentclass[
  letterpaper,
  DIV=11,
  numbers=noendperiod]{scrartcl}
\usepackage{xcolor}
\usepackage{amsmath,amssymb}
\setcounter{secnumdepth}{-\maxdimen} % remove section numbering
\usepackage{iftex}
\ifPDFTeX
  \usepackage[T1]{fontenc}
  \usepackage[utf8]{inputenc}
  \usepackage{textcomp} % provide euro and other symbols
\else % if luatex or xetex
  \usepackage{unicode-math} % this also loads fontspec
  \defaultfontfeatures{Scale=MatchLowercase}
  \defaultfontfeatures[\rmfamily]{Ligatures=TeX,Scale=1}
\fi
\usepackage{lmodern}
\ifPDFTeX\else
  % xetex/luatex font selection
\fi
% Use upquote if available, for straight quotes in verbatim environments
\IfFileExists{upquote.sty}{\usepackage{upquote}}{}
\IfFileExists{microtype.sty}{% use microtype if available
  \usepackage[]{microtype}
  \UseMicrotypeSet[protrusion]{basicmath} % disable protrusion for tt fonts
}{}
\makeatletter
\@ifundefined{KOMAClassName}{% if non-KOMA class
  \IfFileExists{parskip.sty}{%
    \usepackage{parskip}
  }{% else
    \setlength{\parindent}{0pt}
    \setlength{\parskip}{6pt plus 2pt minus 1pt}}
}{% if KOMA class
  \KOMAoptions{parskip=half}}
\makeatother
% Make \paragraph and \subparagraph free-standing
\makeatletter
\ifx\paragraph\undefined\else
  \let\oldparagraph\paragraph
  \renewcommand{\paragraph}{
    \@ifstar
      \xxxParagraphStar
      \xxxParagraphNoStar
  }
  \newcommand{\xxxParagraphStar}[1]{\oldparagraph*{#1}\mbox{}}
  \newcommand{\xxxParagraphNoStar}[1]{\oldparagraph{#1}\mbox{}}
\fi
\ifx\subparagraph\undefined\else
  \let\oldsubparagraph\subparagraph
  \renewcommand{\subparagraph}{
    \@ifstar
      \xxxSubParagraphStar
      \xxxSubParagraphNoStar
  }
  \newcommand{\xxxSubParagraphStar}[1]{\oldsubparagraph*{#1}\mbox{}}
  \newcommand{\xxxSubParagraphNoStar}[1]{\oldsubparagraph{#1}\mbox{}}
\fi
\makeatother

\usepackage{color}
\usepackage{fancyvrb}
\newcommand{\VerbBar}{|}
\newcommand{\VERB}{\Verb[commandchars=\\\{\}]}
\DefineVerbatimEnvironment{Highlighting}{Verbatim}{commandchars=\\\{\}}
% Add ',fontsize=\small' for more characters per line
\usepackage{framed}
\definecolor{shadecolor}{RGB}{241,243,245}
\newenvironment{Shaded}{\begin{snugshade}}{\end{snugshade}}
\newcommand{\AlertTok}[1]{\textcolor[rgb]{0.68,0.00,0.00}{#1}}
\newcommand{\AnnotationTok}[1]{\textcolor[rgb]{0.37,0.37,0.37}{#1}}
\newcommand{\AttributeTok}[1]{\textcolor[rgb]{0.40,0.45,0.13}{#1}}
\newcommand{\BaseNTok}[1]{\textcolor[rgb]{0.68,0.00,0.00}{#1}}
\newcommand{\BuiltInTok}[1]{\textcolor[rgb]{0.00,0.23,0.31}{#1}}
\newcommand{\CharTok}[1]{\textcolor[rgb]{0.13,0.47,0.30}{#1}}
\newcommand{\CommentTok}[1]{\textcolor[rgb]{0.37,0.37,0.37}{#1}}
\newcommand{\CommentVarTok}[1]{\textcolor[rgb]{0.37,0.37,0.37}{\textit{#1}}}
\newcommand{\ConstantTok}[1]{\textcolor[rgb]{0.56,0.35,0.01}{#1}}
\newcommand{\ControlFlowTok}[1]{\textcolor[rgb]{0.00,0.23,0.31}{\textbf{#1}}}
\newcommand{\DataTypeTok}[1]{\textcolor[rgb]{0.68,0.00,0.00}{#1}}
\newcommand{\DecValTok}[1]{\textcolor[rgb]{0.68,0.00,0.00}{#1}}
\newcommand{\DocumentationTok}[1]{\textcolor[rgb]{0.37,0.37,0.37}{\textit{#1}}}
\newcommand{\ErrorTok}[1]{\textcolor[rgb]{0.68,0.00,0.00}{#1}}
\newcommand{\ExtensionTok}[1]{\textcolor[rgb]{0.00,0.23,0.31}{#1}}
\newcommand{\FloatTok}[1]{\textcolor[rgb]{0.68,0.00,0.00}{#1}}
\newcommand{\FunctionTok}[1]{\textcolor[rgb]{0.28,0.35,0.67}{#1}}
\newcommand{\ImportTok}[1]{\textcolor[rgb]{0.00,0.46,0.62}{#1}}
\newcommand{\InformationTok}[1]{\textcolor[rgb]{0.37,0.37,0.37}{#1}}
\newcommand{\KeywordTok}[1]{\textcolor[rgb]{0.00,0.23,0.31}{\textbf{#1}}}
\newcommand{\NormalTok}[1]{\textcolor[rgb]{0.00,0.23,0.31}{#1}}
\newcommand{\OperatorTok}[1]{\textcolor[rgb]{0.37,0.37,0.37}{#1}}
\newcommand{\OtherTok}[1]{\textcolor[rgb]{0.00,0.23,0.31}{#1}}
\newcommand{\PreprocessorTok}[1]{\textcolor[rgb]{0.68,0.00,0.00}{#1}}
\newcommand{\RegionMarkerTok}[1]{\textcolor[rgb]{0.00,0.23,0.31}{#1}}
\newcommand{\SpecialCharTok}[1]{\textcolor[rgb]{0.37,0.37,0.37}{#1}}
\newcommand{\SpecialStringTok}[1]{\textcolor[rgb]{0.13,0.47,0.30}{#1}}
\newcommand{\StringTok}[1]{\textcolor[rgb]{0.13,0.47,0.30}{#1}}
\newcommand{\VariableTok}[1]{\textcolor[rgb]{0.07,0.07,0.07}{#1}}
\newcommand{\VerbatimStringTok}[1]{\textcolor[rgb]{0.13,0.47,0.30}{#1}}
\newcommand{\WarningTok}[1]{\textcolor[rgb]{0.37,0.37,0.37}{\textit{#1}}}

\usepackage{longtable,booktabs,array}
\newcounter{none} % for unnumbered tables
\usepackage{calc} % for calculating minipage widths
% Correct order of tables after \paragraph or \subparagraph
\usepackage{etoolbox}
\makeatletter
\patchcmd\longtable{\par}{\if@noskipsec\mbox{}\fi\par}{}{}
\makeatother
% Allow footnotes in longtable head/foot
\IfFileExists{footnotehyper.sty}{\usepackage{footnotehyper}}{\usepackage{footnote}}
\makesavenoteenv{longtable}
\usepackage{graphicx}
\makeatletter
\newsavebox\pandoc@box
\newcommand*\pandocbounded[1]{% scales image to fit in text height/width
  \sbox\pandoc@box{#1}%
  \Gscale@div\@tempa{\textheight}{\dimexpr\ht\pandoc@box+\dp\pandoc@box\relax}%
  \Gscale@div\@tempb{\linewidth}{\wd\pandoc@box}%
  \ifdim\@tempb\p@<\@tempa\p@\let\@tempa\@tempb\fi% select the smaller of both
  \ifdim\@tempa\p@<\p@\scalebox{\@tempa}{\usebox\pandoc@box}%
  \else\usebox{\pandoc@box}%
  \fi%
}
% Set default figure placement to htbp
\def\fps@figure{htbp}
\makeatother





\setlength{\emergencystretch}{3em} % prevent overfull lines

\providecommand{\tightlist}{%
  \setlength{\itemsep}{0pt}\setlength{\parskip}{0pt}}



 


\usepackage{fvextra}
\DefineVerbatimEnvironment{Highlighting}{Verbatim}{breaklines,commandchars=\\\{\}}\DefineVerbatimEnvironment{verbatim}{Verbatim}{breaklines,commandchars=\\\{\}}
\KOMAoption{captions}{tableheading}
\makeatletter
\@ifpackageloaded{caption}{}{\usepackage{caption}}
\AtBeginDocument{%
\ifdefined\contentsname
  \renewcommand*\contentsname{Table of contents}
\else
  \newcommand\contentsname{Table of contents}
\fi
\ifdefined\listfigurename
  \renewcommand*\listfigurename{List of Figures}
\else
  \newcommand\listfigurename{List of Figures}
\fi
\ifdefined\listtablename
  \renewcommand*\listtablename{List of Tables}
\else
  \newcommand\listtablename{List of Tables}
\fi
\ifdefined\figurename
  \renewcommand*\figurename{Figure}
\else
  \newcommand\figurename{Figure}
\fi
\ifdefined\tablename
  \renewcommand*\tablename{Table}
\else
  \newcommand\tablename{Table}
\fi
}
\@ifpackageloaded{float}{}{\usepackage{float}}
\floatstyle{ruled}
\@ifundefined{c@chapter}{\newfloat{codelisting}{h}{lop}}{\newfloat{codelisting}{h}{lop}[chapter]}
\floatname{codelisting}{Listing}
\newcommand*\listoflistings{\listof{codelisting}{List of Listings}}
\makeatother
\makeatletter
\makeatother
\makeatletter
\@ifpackageloaded{caption}{}{\usepackage{caption}}
\@ifpackageloaded{subcaption}{}{\usepackage{subcaption}}
\makeatother
\usepackage{bookmark}
\IfFileExists{xurl.sty}{\usepackage{xurl}}{} % add URL line breaks if available
\urlstyle{same}
\hypersetup{
  pdftitle={ENVS 193DS Final},
  pdfauthor={Steven Le},
  colorlinks=true,
  linkcolor={blue},
  filecolor={Maroon},
  citecolor={Blue},
  urlcolor={Blue},
  pdfcreator={LaTeX via pandoc}}


\title{ENVS 193DS Final}
\author{Steven Le}
\date{2026-06-04}
\begin{document}
\maketitle

\renewcommand*\contentsname{Table of contents}
{
\setcounter{tocdepth}{3}
\tableofcontents
}

\href{https://github.com/lesteven24/ENVS-193DS_spring-2026_final.git}{GitHub
repo link}

\section{Homework set up}\label{homework-set-up}

\begin{Shaded}
\begin{Highlighting}[]
\CommentTok{\#read in packages}
\FunctionTok{library}\NormalTok{(tidyverse)}
\FunctionTok{library}\NormalTok{(here)}
\FunctionTok{library}\NormalTok{(janitor)}
\FunctionTok{library}\NormalTok{(readxl)}
\FunctionTok{library}\NormalTok{(ggplot2)}
\FunctionTok{library}\NormalTok{(ggeffects)}
\FunctionTok{library}\NormalTok{(flextable)}
\FunctionTok{library}\NormalTok{(here)}
\FunctionTok{library}\NormalTok{(MuMIn)}
\FunctionTok{library}\NormalTok{(DHARMa)}
\FunctionTok{library}\NormalTok{(rstatix)}
\CommentTok{\#stored nest data as an object called nests}
\NormalTok{nests }\OtherTok{\textless{}{-}} \FunctionTok{read\_csv}\NormalTok{(}\FunctionTok{here}\NormalTok{(}\StringTok{"data"}\NormalTok{, }\StringTok{"nest\_data\_final.csv"}\NormalTok{))}
\end{Highlighting}
\end{Shaded}

\section{Problem 1. Research writing}\label{problem-1.-research-writing}

\subsection{a. Transparent statistical
methods}\label{a.-transparent-statistical-methods}

In part 1, they used a logistic regression, which is a generalized
linear model with a binomial error distribution. In part 2, they used a
chi-square test, which is used to identify if there is a relationship
between two categorical variables.

\subsection{b. Suggestions for
rewriting}\label{b.-suggestions-for-rewriting}

\subsubsection{Part 1.}\label{part-1.}

The distance to water edge does seem to predict the probability of
American Avocet microhabitat use.

\subsubsection{Part 2.}\label{part-2.}

There does seem to be a relationship between bird species (American
Avocet, Forster's Tern, Black-necked Stilt) and habitat (managed pond,
non-tidal marsh, salt pond, tidal marsh, or tidal mudflat).

\subsection{c.~Further analysis}\label{c.-further-analysis}

One additional variable that could influence the probability of American
Avocet microhabitat use is the type of behavior they were displaying,
which is a categorical variable. For the American Avocets, the
researchers put behavior into four main categories: feeding (searching,
probing, and dunking bills into the water), roosting (standing, sitting,
preening), breeding (nest-building, incubating, brooding chicks, alarm
calling, breeding display, copulation, or predator distraction display),
and other (walking, swimming, or flushing). This variable might
influence American Avocet habitat use as different behaviors can change
the duration of how long they remain in a particular habitat.

\section{Problem 2. Data analysis}\label{problem-2.-data-analysis}

\subsection{a. Clean and wrangle}\label{a.-clean-and-wrangle}

\begin{Shaded}
\begin{Highlighting}[]
\CommentTok{\#created new object from nests}
\NormalTok{nests\_clean }\OtherTok{\textless{}{-}}\NormalTok{ nests }\SpecialCharTok{|\textgreater{}}
\CommentTok{\#cleaned column names    }
  \FunctionTok{clean\_names}\NormalTok{() }\SpecialCharTok{|\textgreater{}} 
\CommentTok{\#selected for tree height, ant\_species, and case\_control columns }
  \FunctionTok{select}\NormalTok{(case\_control, ant\_species, height\_cm) }\SpecialCharTok{|\textgreater{}} 
\CommentTok{\#reordered names in ant\_species column }
  \FunctionTok{mutate}\NormalTok{(}\AttributeTok{ant\_species =} \FunctionTok{as\_factor}\NormalTok{(ant\_species),}
         \AttributeTok{ant\_species =} \FunctionTok{fct\_relevel}\NormalTok{(ant\_species,}
                                   \StringTok{"AB"}\NormalTok{,}
                                   \StringTok{"RRB"}\NormalTok{,}
                                   \StringTok{"BBR"}\NormalTok{)) }\SpecialCharTok{|\textgreater{}} 
\CommentTok{\#filtered ant\_species column to only have three required species  }
  \FunctionTok{filter}\NormalTok{(ant\_species }\SpecialCharTok{\%in\%} \FunctionTok{c}\NormalTok{(}\StringTok{"AB"}\NormalTok{, }\StringTok{"RRB"}\NormalTok{, }\StringTok{"BBR"}\NormalTok{)) }\SpecialCharTok{|\textgreater{}} 
\CommentTok{\#created new column with full scientific name for each ant species }
  \FunctionTok{mutate}\NormalTok{(}\AttributeTok{scientific\_name =} \FunctionTok{case\_match}\NormalTok{(ant\_species,}
    \StringTok{"AB"} \SpecialCharTok{\textasciitilde{}} \StringTok{"Crematogaster sjostedti"}\NormalTok{,}
    \StringTok{"RRB"} \SpecialCharTok{\textasciitilde{}} \StringTok{"Crematogaster mimosae"}\NormalTok{,}
    \StringTok{"BBR"} \SpecialCharTok{\textasciitilde{}} \StringTok{"Crematogaster nigriceps"}\NormalTok{)) }
\CommentTok{\#displayed structure of nests\_clean data frame}
\FunctionTok{str}\NormalTok{(nests\_clean)}
\end{Highlighting}
\end{Shaded}

\begin{verbatim}
tibble [270 x 4] (S3: tbl_df/tbl/data.frame)
 $ case_control   : num [1:270] 1 0 0 0 0 1 0 0 0 1 ...
 $ ant_species    : Factor w/ 5 levels "AB","RRB","BBR",..: 2 2 2 2 1 2 3 2 1 3 ...
 $ height_cm      : num [1:270] 392 310 513 154 324 ...
 $ scientific_name: chr [1:270] "Crematogaster mimosae" "Crematogaster mimosae" "Crematogaster mimosae" "Crematogaster mimosae" ...
\end{verbatim}

\begin{Shaded}
\begin{Highlighting}[]
\CommentTok{\#displayed 10 random rows from the data frame but made it separate from all codes above by using separate pipe operator }
\NormalTok{nests\_clean }\SpecialCharTok{|\textgreater{}}  \FunctionTok{slice\_sample}\NormalTok{(}\AttributeTok{n =} \DecValTok{10}\NormalTok{)}
\end{Highlighting}
\end{Shaded}

\begin{verbatim}
# A tibble: 10 x 4
   case_control ant_species height_cm scientific_name        
          <dbl> <fct>           <dbl> <chr>                  
 1            0 RRB             200.  Crematogaster mimosae  
 2            1 RRB             512.  Crematogaster mimosae  
 3            0 RRB              99   Crematogaster mimosae  
 4            0 AB              132.  Crematogaster sjostedti
 5            0 BBR             140   Crematogaster nigriceps
 6            0 RRB             368   Crematogaster mimosae  
 7            0 RRB              76.5 Crematogaster mimosae  
 8            0 RRB             165   Crematogaster mimosae  
 9            1 BBR             246   Crematogaster nigriceps
10            0 RRB             132   Crematogaster mimosae  
\end{verbatim}

\subsection{b. Response variable}\label{b.-response-variable}

For the case\_control column, the 1s mean that a tree does contain a
nest and the 0s mean that a tree does not contain a nest (is an
``available'' tree for comparison).

\subsection{c.~Purpose of study}\label{c.-purpose-of-study}

Tree height is a continuous (numeric) variable and could influence the
probability of nest occurrence as if tree height were to increase, the
probability of nest occurrence would increase because there would most
likely be less predators higher up on trees (Results section, paragraph
2). Acacia ant species is a categorical (nominal variable) and could
influence the probability of nest occurrence as if a particular ant
species at a nesting site was more aggressive, the probability of nest
occurrence could increase as it would reduce nest predation risks
(Discussion section, paragraph 1).

\subsection{d.~Table of models}\label{d.-table-of-models}

{\def\LTcaptype{none} % do not increment counter
\begin{longtable}[]{@{}
  >{\raggedright\arraybackslash}p{(\linewidth - 6\tabcolsep) * \real{0.1728}}
  >{\centering\arraybackslash}p{(\linewidth - 6\tabcolsep) * \real{0.1605}}
  >{\centering\arraybackslash}p{(\linewidth - 6\tabcolsep) * \real{0.1605}}
  >{\centering\arraybackslash}p{(\linewidth - 6\tabcolsep) * \real{0.5062}}@{}}
\toprule\noalign{}
\begin{minipage}[b]{\linewidth}\raggedright
Model number
\end{minipage} & \begin{minipage}[b]{\linewidth}\centering
Tree height
\end{minipage} & \begin{minipage}[b]{\linewidth}\centering
Ant species
\end{minipage} & \begin{minipage}[b]{\linewidth}\centering
Model description
\end{minipage} \\
\midrule\noalign{}
\endhead
\bottomrule\noalign{}
\endlastfoot
0 & X & X & All predictors (no interaction term, +) \\
1 & X & X & All predictors (interaction term, *) \\
\end{longtable}
}

\subsection{e. Run the models}\label{e.-run-the-models}

\begin{Shaded}
\begin{Highlighting}[]
\CommentTok{\#model 0: all predictors, no interaction term}
\NormalTok{model0 }\OtherTok{\textless{}{-}} \FunctionTok{glm}\NormalTok{(case\_control }\SpecialCharTok{\textasciitilde{}}\NormalTok{ height\_cm }\SpecialCharTok{+}\NormalTok{ scientific\_name,}
  \AttributeTok{data =}\NormalTok{ nests\_clean,}
\CommentTok{\#response variable is binary so test should be logistic regression}
  \AttributeTok{family =}\NormalTok{ binomial)}

\CommentTok{\#model 1: all predictors, interaction term}
\NormalTok{model1 }\OtherTok{\textless{}{-}} \FunctionTok{glm}\NormalTok{(case\_control }\SpecialCharTok{\textasciitilde{}}\NormalTok{ height\_cm }\SpecialCharTok{*}\NormalTok{ scientific\_name,}
  \AttributeTok{data =}\NormalTok{ nests\_clean,}
\CommentTok{\#response variable is binary so test should be logistic regression}
  \AttributeTok{family =}\NormalTok{ binomial)}
\end{Highlighting}
\end{Shaded}

\subsection{f.~Select the best model}\label{f.-select-the-best-model}

\begin{Shaded}
\begin{Highlighting}[]
\CommentTok{\#calculated AIC for each model}
\FunctionTok{AICc}\NormalTok{(}
\NormalTok{  model0,}
\NormalTok{  model1) }\SpecialCharTok{|\textgreater{}} 
\CommentTok{\#arranged output in order of AIC}
  \FunctionTok{arrange}\NormalTok{(AICc)}
\end{Highlighting}
\end{Shaded}

\begin{verbatim}
       df     AICc
model1  6 238.1236
model0  4 258.8940
\end{verbatim}

The best model as determined by Akaike's Information Criterion (AIC) was
the model where tree height (cm) and acacia ant species were interacting
terms for the probability of nest occurrence.

\subsection{g. Check the diagnostics}\label{g.-check-the-diagnostics}

\subsection{h. Visualize the model
predictions}\label{h.-visualize-the-model-predictions}

\subsection{i. Write a caption for your
figure}\label{i.-write-a-caption-for-your-figure}

\subsection{j. Calculate model
predictions}\label{j.-calculate-model-predictions}

\subsection{k. Interpret your results}\label{k.-interpret-your-results}

\section{Problem 3. Affective and exploratory
visualizations}\label{problem-3.-affective-and-exploratory-visualizations}

\subsection{a. Comparing
visualizations}\label{a.-comparing-visualizations}

\subsection{b. Designing an analysis}\label{b.-designing-an-analysis}

\subsection{c.~Check any assumptions and run your
analysis}\label{c.-check-any-assumptions-and-run-your-analysis}

\subsection{d.~Create a visualization}\label{d.-create-a-visualization}

\subsection{e. Write a caption}\label{e.-write-a-caption}

\subsection{f.~Write about your
results}\label{f.-write-about-your-results}

\subsection{g. Sharing your affective
visualization}\label{g.-sharing-your-affective-visualization}

COMPLETED ALREADY




\end{document}
